\section{Related work}
\label{sec:related}

\paragraph{Memorisation in genomic models, and tokenisation in text.}
\cite{nemecek2026quantifying} is the direct predecessor: four models on four datasets under three
risk scores, showing that they differ and attributing it to architecture. Its most striking result is
also the one that motivates ours --- Evo, a 7B byte-level model adapted with LoRA, recovers
\NUM{100\%} of planted canaries on all three real genomes even at a single insertion, where
DNABERT-2 resists extraction entirely --- because the two differ in scale, tokenisation,
architecture, objective and fine-tuning strategy at once, so the comparison cannot say which of
those is responsible. We ask which design choice is, which requires holding everything else fixed
rather than comparing releases. We reuse its canary construction and two of its datasets so the
numbers can be read against each other, our controlled backbone is its 12.9M character-level model,
and we audit a larger panel of released checkpoints in Sec.~\ref{sec:published}. In text,
\cite{kharitonov2021bpe} showed that a larger BPE vocabulary makes a model reproduce more of its
training data and conjectured that the shorter token sequences are the cause --- untestable there,
because re-tokenising changes length by construction. A spaced $k$-mer tokeniser pads each content
token with unscored filler, so the vocabulary changes while every window stays 288 slots long; the
coarse-tokeniser behaviour follows the vocabulary, not the length (Appendix~A3). Our contribution is
not that vocabulary matters but the mechanism, corpus-level token rarity, and the finding that the
objective dominates the tokeniser.

\paragraph{Measuring memorisation, and scoring masked models.}
Our three rulers are standard, adapted to a four-letter alphabet with a known chance level: surprise
per nucleotide is the per-token loss of \cite{carlini2023quantifying} divided by the nucleotides a
token covers, which is what makes tokenisers comparable; extraction with a revealed prefix is their
discoverable-extraction protocol; recognition against a candidate pool is canary exposure in the
sense of \cite{jagielski2023canary}, which warns against reading exposure as a probability, so we
report the rank-1 rate directly. \cite{tirumala2022memorization} shows memorisation preceding
overfitting, which is what our matched-held-out-loss comparison controls for, and
\cite{morris2025capacity} supplies the capacity estimate we use. Masked models have no left-to-right
likelihood, so we score them by pseudo-log-likelihood \cite{salazar2020pll} and by a
prefix-conditioned scorer, keeping both because they disagree at low mask rates and calibrating each
against the \SI{2.000}{\bit\per\nucleotide} chance level (Appendix~A1). \cite{mireshghallah2022mlm}
attacks masked models with a likelihood-ratio test; our recognition ruler is deliberately weaker,
and Sec.~\ref{sec:human} shows that this is what matters --- on real haplotypes the $k$-mer
checkpoints look safe under ranking while reproducing records verbatim.

\paragraph{Infilling objectives and infilling attacks.}
Three concurrent results in text bear on our objective axis, none of them genomic.
\cite{vonarx2026dynamics} is the closest analogue of our central comparison: matched Llama 3.2
models pretrained with fill-in-the-middle against left-to-right, finding that infilling recovers
short or partial spans more often while left-to-right assigns high confidence to long exact
continuations --- the direction we measure on DNA --- and that extraction under infilling grows
roughly linearly with repetitions, which is why we report a repetition tier.
\cite{luo2026diffusion} derives the generation-side mechanism: extraction probability increases
monotonically with sampling resolution, with autoregressive decoding as the maximal-resolution
limit. \cite{wang2026infilling} changes how we measure: prefix-only probing understates extraction
from models that denoise arbitrary positions, by up to a factor of three on LLaDA-8B and Dream-7B.
Our extraction ruler is prefix-conditioned, so we re-ask every masked cell in its own idiom and
report both protocols (Sec.~\ref{sec:mitigation}).
